% =====================================================================
% constraints_standalone.tex -- compilable wrapper around the four
% fragments, so they can be proof-read on their own:
%
%     pdflatex constraints_standalone.tex
%
% To use them in a talk instead, \input the fragments directly into
% beamer frames; none of them defines a preamble or a sectioning level
% above \paragraph.
% =====================================================================
\documentclass[11pt]{article}
\usepackage[margin=1in]{geometry}
\usepackage{amsmath, amssymb}
\usepackage{textcomp}
\usepackage[T1]{fontenc}

\title{Interface-width and mesh constraints in \texttt{comp\_eps.py}}
\author{enceladus\_DSM}
\date{}

\begin{document}
\maketitle

\noindent
Every equation below is a transcription of
\texttt{preprocess/comp\_eps.py}, which is the single source of truth for
$\epsilon$, the mesh, and the derived phase-field coefficients. References:
Kaempfer \& Plapp, \emph{Phys.\ Rev.\ E} \textbf{79}, 031502 (2009)
[``K\&P''] and Moure \& Fu, \emph{Cryst.\ Growth Des.} \textbf{24}, 7808
(2024) [``M\&F''].

\section{Kinetics: the inputs every bound is evaluated against}
% =====================================================================
% constraints_kinetics.tex -- the state -> alpha_c -> beta chain that
% every eps bound is evaluated against.
%
% Source of truth: preprocess/comp_eps.py (rho_vs_sat, Dv_T, beta_HK,
% capillary_length, sigma0, alpha_libbrecht).  Fragment: no preamble,
% \input it into a beamer frame or a report.  Requires amsmath.
% =====================================================================

\paragraph{Thermodynamic state at the run temperature $T$.}
Everything below is evaluated once, at the single temperature the mesh is
sized for; the solver aborts if the run temperature differs from it by more
than $1\,^\circ$C (\texttt{-eps\_valid\_temp}).
%
\begin{align}
  \rho_{vs}(T) &= \rho_{a}\,b\,\frac{p_{vs}(T)}{p_{atm}-p_{vs}(T)},
  &&\text{saturation vapour density (ASHRAE)} \label{eq:rhovs}\\[2pt]
  D_v(T) &= D_{v0}\left(\frac{T}{273.15\,\mathrm{K}}\right)^{1.81},
  &&D_{v0} = 2.178\times10^{-5}\ \mathrm{m^2/s} \label{eq:Dv}\\[2pt]
  d_0(T) &= \frac{\gamma\,(m/\rho_i)}{k_B T} \;=\; \frac{\gamma V_m}{R T},
  &&\text{capillary length (K\&P Eq.\ 13)} \label{eq:d0}\\[2pt]
  D_i = \frac{\kappa_i}{C_i}, \quad D_a &= \frac{\kappa_a}{C_a}, \quad
  D^{*}_{ia} = \tfrac{1}{2}\left(D_i + D_a\right),
  &&\text{thermal channels} \label{eq:Dheat}
\end{align}

\paragraph{Attachment kinetics.}
$\alpha_c$ is the model's only genuinely free parameter. It fixes the
Hertz--Knudsen coefficient, and through it every bound in
\texttt{constraints\_eps.tex}:
%
\begin{align}
  \beta_{HK}(T,\alpha_c) &= \frac{1}{\alpha_c}\sqrt{\frac{2\pi m}{k_B T}},
  &&\text{\emph{scaled} form (K\&P $\beta'$)} \label{eq:betahk}\\[2pt]
  \beta_{sub}(T,\alpha_c) &= \frac{\beta_{HK}}{\rho_{vs}(T)/\rho_i}
  \;=\; \frac{\rho_i}{\rho_{vs}(T)}\,\frac{1}{\alpha_c}
        \sqrt{\frac{2\pi m}{k_B T}},
  &&\text{\emph{unscaled} (K\&P $\beta_0 \equiv$ M\&F $\beta_{sub}$)}
  \label{eq:betasub}
\end{align}
%
The two differ by the factor $\rho_{vs}/\rho_i \sim 10^{-6}$, so mixing them
is a six-order-of-magnitude error. Eqs.~\eqref{eq:bheat} and \eqref{eq:bvapor}
below take the scaled $\beta_{HK}$; Eq.~\eqref{eq:bkinetic} takes the unscaled
$\beta_{sub}$. The solver is passed \texttt{-beta\_sub0} $=\beta_{sub}$ and
\texttt{-d0\_sub0} $=d_0$ and rescales both internally.

\paragraph{Where $\alpha_c$ comes from.}
Two models are implemented. Libbrecht's nucleation form makes $\alpha_c$ a
function of the local state,
%
\begin{equation}
  \alpha_c(T,\sigma) \;=\; \exp\!\left[-\,\frac{\sigma_0(T)}{\sigma}\right],
  \qquad
  \sigma \;=\; \frac{\lvert \rho_v - \rho_{vs}(T)\rvert}{\rho_{vs}(T)},
  \label{eq:libbrecht}
\end{equation}
%
with $\sigma_0(T)$ a log--log interpolant of the Libbrecht (2017) table in
$\lvert T \rvert$, held flat beyond both ends, and a hard floor
$\alpha_c = 10^{-30}$ whenever $\sigma < \sigma_0/\ln(10^{30})$.
Section \texttt{constraints\_libbrecht.tex} shows why
Eq.~\eqref{eq:libbrecht} is not usable here. The fallback is a constant
$\alpha_c$ inside the literature band $10^{-3} \le \alpha_c \le 10^{-1}$
(Libbrecht 2017; Braun, Fourteau \& L\"owe 2024).


\section{The bounds on $\epsilon$, and the mesh}
% =====================================================================
% constraints_eps.tex -- the four bounds on the interface width eps and
% the mesh rule that follows from it.
%
% Source of truth: preprocess/comp_eps.py, compute_eps().  Fragment: no
% preamble; requires amsmath.  Read constraints_kinetics.tex first --
% it defines beta_HK, beta_sub, d0, D_v and D*_ia.
% =====================================================================

\paragraph{Width convention.}
This solver uses Moure \& Fu's well $\tfrac12\phi^2(1-\phi)^2$ with an
$\epsilon^2$ gradient term, so the equilibrium profile is logistic,
%
\begin{equation}
  \phi(x) \;=\; \frac{1}{1+e^{-x/\epsilon}}
          \;=\; \tfrac12\left[1+\tanh\!\left(\frac{x}{2\epsilon}\right)\right],
  \label{eq:profile}
\end{equation}
%
and $\epsilon$ is a \emph{decay length}, not the visible band width: the
$5\%$--$95\%$ transition spans $\approx 6\epsilon$ and the $1\%$--$99\%$
transition $\approx 9.2\epsilon$. Kaempfer \& Plapp's $\pm1$ well has
$W = \sqrt{2}\,\epsilon$; the asymptotic constants used below,
$a_1 = 5$ and $a_2 = 0.1581$, are the ones for \emph{this} well.

\paragraph{The four bounds.}
Three are Kaempfer \& Plapp (2009) sharp-interface conditions and carry a
common safety factor $s$ (default $s = \tfrac12$); the fourth is geometric
and does not.
%
\begin{align}
  \textbf{B-HEAT:}    &&  \epsilon &\;\le\; s\,D_{heat}\,\beta_{HK}
    && \text{K\&P Eq.\ (43a/b)} \label{eq:bheat}\\[2pt]
  \textbf{B-VAPOR:}   &&  \epsilon &\;\le\; s\,D_v\,\beta_{HK}
    && \text{K\&P Eq.\ (43c)} \label{eq:bvapor}\\[2pt]
  \textbf{B-KINETIC:} &&  \epsilon &\;\le\; s\,\frac{d_0}{\beta_{sub}\,v_n}
    && \text{K\&P Eq.\ (45)} \label{eq:bkinetic}\\[2pt]
  \textbf{B-CURV:}    &&  \epsilon &\;\le\; 0.05\,R_{ave}
    && \text{geometric, K\&P's } W \ll 1/\chi \label{eq:bcurv}
\end{align}
%
and $\epsilon$ is taken as the smallest of the four:
%
\begin{equation}
  \epsilon \;=\; \min\Big\{\,
      s\,\min\big(D_{heat}\beta_{HK},\; D_v\beta_{HK},\;
                  d_0/(\beta_{sub}v_n)\big),\;\; 0.05\,R_{ave}\,\Big\}.
  \label{eq:epsdef}
\end{equation}

\paragraph{What each bound is for.}
\begin{description}
  \item[B-HEAT, B-VAPOR] control the thin-interface expansion parameter
    $\delta = a_1 a_2\,\epsilon/(D\beta_{HK})$, so on whichever of the two
    channels binds, $\delta = a_1 a_2 s = 0.79\,s$ --- that is what the
    ``$\ll$'' in K\&P Eq.\ (43) means quantitatively. $\tau_{sub}$ already
    compensates $\delta$ to first order, so $\delta$ is the size of the
    correction being applied, not the residual; the residual is
    $\mathcal{O}(\delta^2)$.
  \item[$D_{heat}$] is $D^{*}_{ia}$ by default --- the diffusivity the solver
    actually assembles and that $\tau_{sub}$ compensates against. The
    single-sided ice channel $D_i = \kappa_i/C_i$ is $\sim6\times$ tighter and
    is knowingly violated (enforcing it drove one geometry to $\sim10^{9}$
    nodes); M\&F violate it by $100\times$ in print and say so. It is reported
    as the diagnostic $\delta_{ice}$ instead of enforced.
  \item[B-KINETIC] keeps the kinetic undercooling $\beta_{sub}v_n$ small
    against the capillary term $d_0/\epsilon$. It is the \emph{only} bound
    that involves the front velocity $v_n$, and the only one that tightens
    as $\alpha_c$ falls.
  \item[B-CURV] is deliberately crude. It is one to two decades looser than
    the binding bound in every case we run, and the smallest true curvature
    (a fresh neck fillet) starts near-singular and then relaxes, so sizing a
    mesh from it is both unstable and needlessly expensive.
\end{description}

\paragraph{Mesh.}
The element size is fixed to $\epsilon$ by
%
\begin{equation}
  h \;=\; \frac{\epsilon}{\sqrt{2}},
  \qquad
  N_d \;=\; \left\lceil \frac{L_d}{h} \right\rceil
        \;=\; \left\lceil \frac{\sqrt{2}\,L_d}{\epsilon} \right\rceil,
  \qquad d \in \{x,y,z\},
  \label{eq:mesh}
\end{equation}
%
i.e.\ two elements per asymptotic width $W = \sqrt{2}\epsilon$, which is
$8.3$ elements across $\phi = 0.05$--$0.95$ and $13.0$ across
$\phi = 0.01$--$0.99$. Cost therefore scales as $\epsilon^{-d}$: the sizer
flags a configuration as \emph{intractable} once
$\prod_d N_d > 5\times10^{7}$.

\paragraph{Diagnostics that are reported but never enforced.}
\begin{equation}
  \delta_{X} = \frac{a_1 a_2\,\epsilon}{D_X\,\beta_{HK}},
  \qquad
  \mathrm{Pe} = \frac{\epsilon\,v_n}{D_v},
  \qquad
  \frac{\epsilon}{R_{ave}},
  \qquad
  \frac{\beta_{sub}v_n}{d_0\chi},
  \qquad
  \xi_v \;\ge\; 10\,\frac{\rho_{vs}}{\rho_i}.
  \label{eq:diagnostics}
\end{equation}
%
$\epsilon/R_{ave}$ is noted above $5\%$ and warned above $10\%$;
$\beta_{sub}v_n/(d_0\chi)$ says whether the Gibbs--Thomson condition is
kinetics- or capillarity-dominated; the last is the temporal-scaling
condition that keeps the unscaled vapour storage term under ${\sim}10\%$ of
the source.


\section{What is derived from $\epsilon$ afterwards}
% =====================================================================
% constraints_derived.tex -- the phase-field parameters derived from eps
% once the bounds of constraints_eps.tex have fixed it.
%
% Source of truth: preprocess/comp_eps.py, derived_pf_params()
% (Moure & Fu 2024, SI Eq. 9).  Fragment: no preamble; requires amsmath.
% =====================================================================

\paragraph{From $\epsilon$ to the solver's coefficients.}
Once $\epsilon$ is fixed, Moure \& Fu's SI Eq.~(9) inverts the
sharp-interface limit to give the coefficients the solver assembles. With
$\rho_{rat} = \rho_{vs}(T)/\rho_i$:
%
\begin{align}
  d_0\,\rho_{rat} &= \frac{a_1\,\epsilon}{\lambda_{sub}}
  &&\Longrightarrow&
  \lambda_{sub} &= \frac{a_1\,\epsilon}{d_0\,\rho_{rat}}, \label{eq:lam}\\[4pt]
  \beta_{HK} &= a_1\left[\frac{\tau_{sub}}{\epsilon\,\lambda_{sub}}
     \;-\; \frac{a_2\,\epsilon}{D^{*}_{ia}}
     \;-\; \frac{a_2\,\epsilon}{D_v}\right]
  &&\Longrightarrow&
  \tau_{sub} &= \epsilon\,\lambda_{sub}\left[
     \frac{\beta_{HK}}{a_1}
     + \frac{a_2\,\epsilon}{D^{*}_{ia}}
     + \frac{a_2\,\epsilon}{D_v}\right], \label{eq:tau}
\end{align}
%
and then
%
\begin{equation}
  M_{sub} \;=\; \frac{\epsilon}{3\,\tau_{sub}},
  \qquad
  \alpha_{sub} \;=\; \frac{\lambda_{sub}}{\tau_{sub}},
  \qquad\text{so}\qquad
  \frac{\alpha_{sub}}{M_{sub}} \;=\; \frac{3\lambda_{sub}}{\epsilon}
  \ \text{exactly.}
  \label{eq:Msub}
\end{equation}
%
That last identity is why \texttt{-mob\_sub} and \texttt{-alph\_sub} are
never written into an options file: the solver derives both from
$(\epsilon,\ \beta_{sub},\ d_0)$, and overriding either alone breaks the
K\&P consistency ratio.

\paragraph{Two conditions that must hold for Eq.~\eqref{eq:tau} to mean
anything.}
\begin{align}
  \tau_{sub} &> 0
  &&\text{else the asymptotic expansion is ill-posed;} \label{eq:tauplus}\\[2pt]
  f_{kin} \;=\; \frac{\beta_{HK}/a_1}
       {\dfrac{\beta_{HK}}{a_1}
        + \dfrac{a_2\epsilon}{D^{*}_{ia}}
        + \dfrac{a_2\epsilon}{D_v}} &\;\gtrsim\; 0.5
  &&\text{else $\tau_{sub}$ is dominated by thin-interface}
  \label{eq:fkin}
\end{align}
%
corrections rather than by the physical attachment kinetics, and the
sharp-interface analogy is being stretched: $f_{kin} < 0.5$ is flagged as
marginal, $f_{kin} < 0.1$ as a warning.

\paragraph{Physical constants (identical to \texttt{material\_properties.c}).}
\begin{center}\small
\begin{tabular}{llll}
\hline
symbol & value & symbol & value \\
\hline
$m$ (H$_2$O molecule) & $2.9915\times10^{-26}$ kg &
$\rho_i$              & $919$ kg/m$^3$ \\
$k_B$                 & $1.380649\times10^{-23}$ J/K &
$\gamma$ (ice--vapour)& $0.109$ J/m$^2$ \\
$\kappa_i$            & $2.29$ W/m/K &
$C_i$                 & $1.8\times10^{6}$ J/m$^3$/K \\
$\kappa_a$            & $0.02$ W/m/K &
$C_a$                 & $1.341\times1004$ J/m$^3$/K \\
$D_{v0}$              & $2.178\times10^{-5}$ m$^2$/s &
$a_1,\ a_2$           & $5,\ 0.1581$ \\
\hline
\end{tabular}
\end{center}


\section{Libbrecht (2017) $\sigma_0(T)$ pushed through the chain}
% =====================================================================
% constraints_libbrecht.tex -- what happens when alpha_c is taken from
% Libbrecht (2017) sigma0(T) and pushed through constraints_eps.tex.
%
% Numbers regenerated by:
%   python preprocess/plot_libbrecht_constraints.py \
%       --out studies/libbrecht_kinetics
% Fragment: no preamble; requires amsmath, booktabs (or drop \toprule etc.).
% =====================================================================

\paragraph{The chain.}
Substituting Eq.~\eqref{eq:libbrecht} into Eq.~\eqref{eq:betasub} and then
into the kinetic bound Eq.~\eqref{eq:bkinetic} and the mesh rule
Eq.~\eqref{eq:mesh} gives a single closed expression:
%
\begin{equation}
  N_x
  \;=\;
  \left\lceil
    \frac{\sqrt{2}\,L_x}{\epsilon}
  \right\rceil,
  \qquad
  \epsilon \;\le\; s\,\frac{d_0\,\rho_{vs}}{\rho_i}
     \left(\frac{k_B T}{2\pi m}\right)^{1/2}
     \frac{1}{v_n}\;
     \exp\!\left[-\frac{\sigma_0(T)}{\sigma}\right]
  \quad\Longrightarrow\quad
  N_x \;\propto\; \exp\!\left[+\frac{\sigma_0(T)}{\sigma}\right].
  \label{eq:chain}
\end{equation}
%
The mesh is exponential in $\sigma_0/\sigma$. Two facts then decide
everything.

\paragraph{1. The supersaturation gap.}
Libbrecht's chamber measurements are taken at
$\sigma \sim 10^{-2}$ to $10^{-1}$. Our problem runs one to two decades
lower: the imposed Molaro wall undersaturation is $1-h = 2.7\times10^{-3}$,
and the Gibbs--Thomson supersaturation at a neck fillet of radius
$\rho_{f}$ is $\sigma = d_0/\rho_{f} \approx 4.5\times10^{-4}$. Under
Eq.~\eqref{eq:libbrecht} that gap is not a two-decade reduction in
$\alpha_c$; it lands in the exponent.

\paragraph{2. The velocity is measured, not assumed.}
Eq.~\eqref{eq:bkinetic} bites only if the front moves. We take $v_n$ from
Molaro et al.\ (2019) Fig.~11 itself --- the integrated neck-growth rate
$v_n = 3.416\times10^{-9}$ m/s. So the statement is not ``pick a velocity
and watch the mesh explode''; it is ``the experiment we are replicating
moves its interface at this speed, and Libbrecht's $\alpha_c$ says that is
unresolvable.''

\paragraph{The result, at $T = -20\,^\circ$C.}
Molaro grain pair, $L_x \times L_y = 450 \times 225\ \mu$m, $s = 1/2$:
%
\begin{center}\footnotesize
\begin{tabular}{lrrrlrr}
\hline
$\sigma$ & $\alpha_c$ & $\beta_{sub}$ [s/m] & $\epsilon$ [m] & binds
  & $N_x$ & 2D nodes \\
\hline
$10^{-1}$ (Libbrecht) & $7.06\times10^{-1}$ & $1.12\times10^{4}$
  & $4.19\times10^{-8}$  & B-HEAT    & $1.5\times10^{4}$  & $1.2\times10^{8}$ \\
$10^{-2}$ (Libbrecht) & $3.07\times10^{-2}$ & $2.58\times10^{5}$
  & $5.77\times10^{-7}$  & B-KINETIC & $1.1\times10^{3}$  & $6.1\times10^{5}$ \\
$2.7\times10^{-3}$ (ours: wall BC) & $2.50\times10^{-6}$ & $3.17\times10^{9}$
  & $4.69\times10^{-11}$ & B-KINETIC & $1.4\times10^{7}$  & $9.2\times10^{13}$ \\
$4.5\times10^{-4}$ (ours: fillet) & $10^{-30}$ (floor) & $7.91\times10^{33}$
  & $1.88\times10^{-35}$ & B-KINETIC & $3.4\times10^{31}$ & $5.7\times10^{62}$ \\
\hline
\multicolumn{7}{l}{\emph{For reference:} production run $\alpha_c = 0.1$,
  $\epsilon = 0.118\ \mu$m, $N_x = 5394$, $14.5$ M nodes.}\\
\multicolumn{7}{l}{\emph{Molecular limit:} $a = (m/\rho_i)^{1/3} = 3.19$ \AA,
  so $L_x/a = 1.4\times10^{6}$ --- one node per water molecule.}\\
\hline
\end{tabular}
\end{center}
%
At our own boundary condition the demand is already $10^{14}$ nodes with
$\epsilon$ half an \AA ngstr\"om; at the fillet supersaturation it is
$10^{62}$ nodes with $\epsilon$ twenty-five orders of magnitude below a water
molecule. This is not an expensive mesh, it is a category error: the
continuum phase field has no meaning below $a$.

\paragraph{Why this is the wrong data, not just an awkward one.}
\begin{enumerate}
  \item \textbf{Libbrecht measured a different experiment.} $\sigma_0(T)$
    comes from single ice crystals grown from vapour in a near-vacuum
    chamber, where the only vapour transport is ballistic and $\sigma$ is
    imposed externally and held one to two decades above ours. Our
    simulations have a continuous vapour field between grains, so $\sigma$ is
    \emph{a solution variable}, set by the local curvature and the wall
    humidity --- not a control knob. Extrapolating a nucleation law two
    decades below where it was calibrated puts the extrapolation entirely
    inside an exponential.
  \item \textbf{The table cannot be fitted anyway.} It is $9$ points, it is
    non-monotonic (a kink at $-6/-7\,^\circ$C that is a digitisation artifact,
    not physics), and $\sigma_0$ spans a factor $27$ over $-2$ to
    $-40\,^\circ$C. Since $\ln\alpha_c = -\sigma_0/\sigma$, no single
    reference $\sigma$ can hold $\alpha_c$ inside one decade across that
    range.
  \item \textbf{The escape route is circular.} Feeding $v_n$ back
    self-consistently from the kinetics, $v_n = d_0/(\beta_{sub}R_{feat})$,
    makes Eq.~\eqref{eq:bkinetic} collapse to $\epsilon \le s\,R_{feat}$ ---
    independent of $\beta_{sub}$, and therefore of $\alpha_c$. The mesh
    becomes affordable, but only because the constraint has been made
    insensitive to the quantity under test. The cost simply moves from
    resolution to runtime: physical $t_{final}$ scales as $1/\alpha_c$.
\end{enumerate}
%
What Libbrecht's data \emph{does} establish is the \emph{sign} of the
temperature dependence and a plausible band for $\alpha_c$. That is how it is
used: as a consistency check on a constant $\alpha_c$ in
$[10^{-3},10^{-1}]$, not as a pointwise law.


\end{document}
